\chapter{Introduzione e Contesto}
\label{cap:introduzione}

% DA COMPLETARE: capitolo introduttivo. Le sezioni 1.1-1.3 richiedono fonti esterne
% (dati Terna e ARERA sulla capacita' rinnovabile installata, documentazione MACSE sulle
% aste per la capacita' di accumulo, letteratura sull'arbitraggio con storage).
% La sezione 1.4 va scritta per ultima, quando gli obiettivi saranno stabilizzati.

\section{La transizione energetica in Italia}
\label{sec:transizione}

% DA COMPLETARE: inquadramento del contesto.
% Elementi da trattare, con dati da reperire e citare:
%   - crescita della capacita' rinnovabile non programmabile (fotovoltaico ed eolico)
%     in Italia negli ultimi anni;
%   - effetto sui prezzi di mercato: compressione dei prezzi diurni, aumento della
%     volatilita' infragiornaliera, comparsa di prezzi negativi;
%   - obiettivi al 2030 del PNIEC e fabbisogno di flessibilita' stimato da Terna.

Il sistema elettrico italiano sta attraversando una fase di trasformazione profonda, guidata
dalla penetrazione crescente di fonti rinnovabili non programmabili. L'effetto di questa
trasformazione sul mercato non riguarda soltanto il livello medio dei prezzi, ma soprattutto
la loro \emph{forma} nell'arco della giornata: la produzione fotovoltaica comprime i prezzi
nelle ore centrali e li lascia elevati nelle ore serali, quando la domanda resta alta e
l'irraggiamento e' venuto meno.

\`E proprio questo divario infragiornaliero a costituire l'opportunit\`a economica per i
sistemi di accumulo, ed \`e la grandezza su cui si concentra il presente lavoro.

\section{I sistemi di accumulo a batteria (BESS)}
\label{sec:bess}

% DA COMPLETARE: descrizione tecnologica ed economica dei BESS.
% Elementi da trattare:
%   - principio di funzionamento e grandezze caratteristiche: potenza nominale (MW),
%     capacita' energetica (MWh), rapporto energia/potenza (durata), efficienza di
%     ciclo completo (round-trip), profondita' di scarica;
%   - servizi che un BESS puo' fornire: arbitraggio sui mercati dell'energia, riserva,
%     regolazione di frequenza, capacity market;
%   - stato dell'arte in Italia: capacita' installata, aste MACSE, prospettive.

\section{Il quesito degli investitori}
\label{sec:quesito}

% DA COMPLETARE: impostazione della domanda di ricerca dal punto di vista dell'investitore.
% Il punto centrale, che il lavoro svolto permette gia' di argomentare, e' il seguente:
% la redditivita' di un impianto di accumulo dipende dal differenziale di prezzo fra ore
% di carica e ore di scarica, ma quel differenziale non e' un dato esogeno: l'accumulo,
% comprando nelle ore basse e vendendo nelle ore alte, riduce esso stesso il differenziale
% da cui trae profitto. Questo effetto di retroazione e' il cuore metodologico della tesi.

\section{Obiettivi della tesi}
\label{sec:obiettivi}

% DA COMPLETARE: da rifinire quando i risultati saranno consolidati.
% Formulazione provvisoria degli obiettivi, coerente con il lavoro svolto finora.

Il lavoro si propone tre obiettivi.

Il primo \`e \textbf{ricostruire il meccanismo di formazione del prezzo} sul Mercato del
Giorno Prima a partire dai dati pubblici delle offerte, in modo da poter calcolare non solo
il prezzo osservato ma anche il prezzo \emph{controfattuale} che si sarebbe formato in
presenza di capacit\`a di accumulo. Questo obiettivo \`e stato raggiunto ed \`e documentato
nel Capitolo~\ref{cap:simulazione}.

Il secondo \`e \textbf{simulare l'effetto di quantit\`a crescenti di accumulo} sui prezzi e
sulle quantit\`a di equilibrio della zona NORD, tenendo conto dell'effetto di retroazione che
l'accumulo stesso esercita sul prezzo.

Il terzo \`e \textbf{determinare il dimensionamento ottimale} dell'accumulo — capacit\`a
energetica, potenza e durata di carica e scarica — e valutarne la redditivit\`a dal punto di
vista di un investitore.

\section{Organizzazione del documento}
\label{sec:organizzazione}

Il Capitolo~\ref{cap:mercato} descrive l'architettura del mercato elettrico italiano e i
meccanismi d'asta rilevanti per il lavoro. Il Capitolo~\ref{cap:strategia} formula il
problema di ottimizzazione della strategia di accumulo. Il Capitolo~\ref{cap:simulazione}
presenta i dati, la ricostruzione delle curve di domanda e offerta e la sua validazione. Il
Capitolo~\ref{cap:valutazione} traduce i risultati in termini economici e finanziari. Il
Capitolo~\ref{cap:conclusioni} riassume i risultati e ne discute i limiti.

\subsection{Citations}
\label{sec:citations}

% DA COMPLETARE: sostituire con le convenzioni effettivamente adottate.

I riferimenti bibliografici sono gestiti con \texttt{biblatex} in stile numerico e sono
raccolti nel file \texttt{bibliografia.bib}. Nel testo una citazione appare come
\texttt{\textbackslash cite\{chiave\}}; le fonti dei dati di mercato sono citate per esteso
alla prima occorrenza.

\subsection{Use of acronyms}
\label{sec:acronimi}

% DA COMPLETARE: valutare se introdurre il pacchetto acronym o glossaries per generare
% automaticamente l'elenco degli acronimi. Per ora l'elenco e' compilato a mano.

Gli acronimi sono definiti per esteso alla prima occorrenza e usati in forma abbreviata nel
seguito. I principali sono riportati nella Tabella~\ref{tab:acronimi}.

\begin{table}[H]
\centering
\caption{Acronimi ricorrenti nel documento.}
\label{tab:acronimi}
\begin{tabular}{@{}ll@{}}
\toprule
\textbf{Sigla} & \textbf{Significato} \\
\midrule
BESS   & \emph{Battery Energy Storage System}, sistema di accumulo a batteria \\
GME    & Gestore dei Mercati Energetici \\
MGP    & Mercato del Giorno Prima \\
MI     & Mercati Infragiornalieri \\
MSD    & Mercato per il Servizio di Dispacciamento \\
MTU    & \emph{Market Time Unit}, periodo elementare di mercato \\
PUN    & Prezzo Unico Nazionale \\
XML    & \emph{eXtensible Markup Language}, formato dei dati pubblicati dal GME \\
\bottomrule
\end{tabular}
\end{table}
